\documentclass[JRDA]{JRDA}

%%%%%%%%%%%%%%%%%%%%%%%%%%%%%%%%%%%%%%%%%%%%%%%%%%%%%%%
% Packages
\usepackage{array}
\usepackage{graphicx}
\usepackage[utf8]{inputenc}

%%%%%%%%%%%%%%%%%%%%%%%%%%%%%%%%%%%%%%%%%%%%%%%%%%%%%%%
% Header
\title{Template for Abstracts for Journée Régionale des Doctorants en Automatique - 2022}
\author[*1]{Presenting Author}
\author[1,2]{Second Author}
\author[2]{Third Author}

\affil[1]{Department name, Institution 1} 
\affil[2]{Department name, Instituion 2} 
% corresponding/presenting author email
\corrauthor{correspondingauthor@email.com}

%%%%%%%%%%%%%%%%%%%%%%%%%%%%%%%%%%%%%%%%%%%%%%%%%%%%%%%
\begin{document}

\maketitle

This document presents the instructions needed to prepare the extended abstract to be included in the proceedings of the JRDA2022. All abstracts should be written in English. In order to achieve graphical consistency, the extended abstract should be preferably written using Microsoft Word. It should start with the paper’s title, authors, authors’ affiliation, address and e-mail, followed by the Abstract's text and the keywords (up to six). Please, indicate the author making the presentation with an asterisk. 

The abstract should be written in MS Word format in single-spaced in 12 point Times New Roman. Do not indent the text paragraphs. It should not exceed 200 words maximum. Please, briefly explain the aim and scope of your study, your materials and methods, and the main conclusions of the study.

 Cite in the text using commands \citep{gargett1984,lilly1983} or \cite{thorpe2005}. 

%%%%%%%%%%%%%%%%%%%%%%%%%%%%%%%%%%%%%%%%%%%%%%%%%%%%%%%
\bibliographystyle{JRDA}
\bibliography{JRDA}

\end{document}